% chapters/05/section-06.tex
\section{Yang Belum Ada pada Gerak Ini}
\label{sec:belum-ada-gerak}

Objek sudah bergerak dengan kecepatan tetap pada setiap laju bingkai. Bagian
penutup ini menjelaskan batas kemampuan program tersebut dan apa yang
ditambahkan pada bab-bab berikutnya.

\subsection{Kecepatan yang Selalu Tetap}

Kecepatan objek pada bab ini hanya berubah karena tombol ditekan atau dilepas.
Selama tombol ditahan, nilainya tetap, dan begitu tombol dilepas, nilainya
langsung menjadi nol. Gerak semacam itu disebut gerak lurus beraturan.

Sebagian besar permainan tidak mencukupi diri dengan gerak semacam itu.
Sebuah benda yang jatuh bergerak makin lama makin cepat; sebuah bola yang
memantul kehilangan sebagian kecepatannya setiap kali menyentuh lantai; sebuah
karakter yang melompat bergerak melambat ketika mendekati titik tertinggi.
Ketiga gejala itu berasal dari \textbf{percepatan}, yaitu perubahan kecepatan
terhadap waktu.

Percepatan ditambahkan pada tahap pembaruan, tepat sebelum kecepatan dipakai
untuk memperbarui posisi. Susunannya menjadi dua tingkat, sama seperti kaidah
Euler pada Bagian~2.4.

\begin{kodekeluaran}
velocity = velocity + acceleration * delta_seconds
position = position + velocity * delta_seconds
\end{kodekeluaran}

Baris pertama itulah yang ditambahkan pada Bab~VI, beserta cara memilih nilai
percepatan agar gravitasi terasa wajar dan lompatan dapat dikendalikan.

\subsection{Objek yang Belum Dapat Bertabrakan}

Objek pada bab ini bergerak menembus apa pun. Ia hanya dibatasi oleh tepi
jendela, dan pembatasan itu pun dilakukan dengan memotong nilai posisi, bukan
dengan memeriksa tabrakan.

Pada Bab~VII, posisi pecahan dan \code{Rect} yang telah Anda selaraskan akan
dipakai untuk memeriksa kapan dua objek bersentuhan. Penyelarasan pada
Bagian~\ref{sec:selaras-rect} justru menjadi bekal untuk
pekerjaan itu: pemeriksaan tabrakan memerlukan \code{Rect} bertipe bilangan
bulat, sedangkan ketelitian gerak memerlukan posisi bertipe pecahan.

\subsection{Objek yang Masih Berupa Variabel Biasa}

Program pada bab ini menggerakkan satu objek yang disimpan pada beberapa
variabel: \code{position}, \code{velocity}, dan \code{player_rect}. Cara itu
memadai untuk satu objek, tetapi menjadi tidak terkelola ketika objeknya
berjumlah sepuluh atau seratus.

Pygame menyediakan kelas \code{Sprite} yang menyatukan gambar, posisi, dan
perilaku sebuah objek menjadi satu kesatuan, serta kelas \code{Group} untuk
mengelola sekumpulan \istilah{sprite} sekaligus. Keduanya dibahas pada
Bab~VII. Sesudah itu, menambahkan objek baru cukup dengan membuat satu objek
lagi, bukan dengan menambah tiga variabel baru untuk setiap objek.

\subsection{Ringkasan Perkembangan}

\begin{table}[htbp]
\centering
\caption{Bagian yang ditambahkan pada gerak objek pada bab-bab berikutnya}
\label{tab:perkembangan-gerak}
{\setlength{\tabcolsep}{4pt}\footnotesize%
\begin{tabularx}{\textwidth}{@{}C{1.6cm} L{5.0cm} Y@{}}
\toprule
\headrow \textbf{Bab} & \textbf{Yang ditambahkan} & \textbf{Akibatnya} \\
\midrule
Bab~VI & Percepatan: gravitasi, impuls lompatan, pantulan &
Kecepatan tidak lagi tetap, melainkan berubah setiap bingkai menurut
percepatan yang bekerja. \\
\addlinespace
Bab~VII & \code{Sprite}, \code{Group}, dan deteksi tabrakan &
Objek menjadi satu kesatuan beserta gambarnya, dan perpindahannya dapat
terhenti karena bertabrakan. \\
\addlinespace
Bab~VIII--XI & Pengendalian NPC &
Kecepatan tidak lagi berasal dari tombol, melainkan dari keputusan yang
diambil program sendiri. \\
\addlinespace
Bab~XIII & Fisika mesin pada Unity &
Perhitungan percepatan dan tabrakan diserahkan kepada mesin, sedangkan kode
Anda mengatur pengaturan dan menanggapi peristiwanya. \\
\bottomrule
\end{tabularx}}
\end{table}

Dengan gambaran itu, tampak bahwa bab ini meletakkan dua hal yang akan dipakai
terus: \textbf{posisi sebagai pecahan} yang diperbarui setiap bingkai, dan
\textbf{perpindahan sebagai hasil kali kecepatan dengan selang waktu}. Kedua
hal itu tidak akan berubah lagi sampai akhir buku, bahkan ketika
perhitungannya diserahkan kepada mesin permainan.
